\documentclass[12pt]{article}
\usepackage{fullpage}
%\usepackage{setspace}
\usepackage{enumitem}
\usepackage{listings,amssymb}
\usepackage{listings}
\usepackage{graphicx,xcolor}
%\doublespacing
\usepackage[T1]{fontenc}
\usepackage{newtxtext,newtxmath}
\lstset{numbers=left, numberstyle=\tiny}
\title{CS 147: Lab 02 \\ Due Date: September 22, 2026}
\author{}
\date{\today}
\begin{document}
\maketitle
%\vspace{-0.5in}
The main objective of this problem set is to become familiar with using the toolchain for the course project.  
Feel free to refer
to the Verilog cheatsheet, Verilog lecture slides linked off of canvas. 

Some advice: Before starting to write any Verilog, you should do the following:
\begin{enumerate}
\item
Break down your design into sub-modules.
\item Define interfaces between these modules.
\item Draw paper and pencil schematics for these modules (these will be handed in as scanned
\texttt{schematic.pdf} file(s)).
\item Then start writing Verilog.
\end{enumerate}

\section*{Problem 1 (20 points)}
Design a 16-bit barrel shifter in Verilog with the following interface. If you need additional information on
barrel shifter design, you can consult the ALU lecture notes. Note that we call this shifter a barrel shifter
because it is capable of shifting and rotating bits all the way around in any direction (like a barrel). You
should \textbf{not} just write a case statement with each constant shift amount – instead think about how the
hardware would be designed for a shifter, and what underlying modules (e.g., muxes) you might be able to
utilize to do shifting. 

\textbf{Inputs:}
\begin{itemize}
\item $[15:0]$ \texttt{InBS} - 16-bit input operand value to be shifted
\item $[3:0]$ \texttt{ShAmt} - 4-bit amount to shift (number of bit positions to shift)
\item $[1:0]$ \texttt{ShiftOper} - shift type, see encoding in table below
\end{itemize}
\textbf{Output:}
\begin{itemize}
\item $[15:0]$ \texttt{OutBS} - 16-bit output operand
\end{itemize}
\begin{center}
\begin{tabular}{|c|c|} \hline 
Opcode & Operation \\ \hline
00 & Rotate left \\ \hline
01 & Shift left \\ \hline
10 & Shift right arithmetic \\ \hline
11 & Shift right logical \\ \hline
\end{tabular}
\end{center}
\textit{(Aside: you should think about if the above 4 opcodes are sufficient to represent all of the shift operations
you need to implement for your project.)}

Verify the design using the testbenches provided by us.  You are also encouraged to write additional tests, as we may use
more tests when grading your assignment.  See class slides from the Verilog lecture on how to write tests.

\section*{Problem 2 [30 points]}

This problem should also be done in Verilog. Design a simple 16-bit ALU. Operations to be performed are
2's Complement ADD, bitwise-OR, bitwise-XOR, bitwise-AND, and the barrel shifter unit from Problem
1. Additionally, it must have the ability to invert either of its data inputs before performing the operation
and have a \texttt{Cin} input (to enable subtraction). Another input line also determines whether the arithmetic to
be performed is signed or unsigned. Use your \texttt{CLA} (from homework 1) in your design, extended as needed
to take things like overflow and sign into account. For all the shift and rotate operations, assume the number
to shift is input \texttt{InA} to your ALU and the shift/rotate amount is bits [3:0] of input \texttt{InB} (\textbf{Note:} \texttt{InA/InB} may
be inverted before reaching the shifter, see below – if either/both is inverted, you should shift/rotate the
inverted values instead).

\begin{center}
\begin{tabular}{|c|c|c|} \hline 
Opcode & Function & Result \\ \hline
000 & rll & Rotate left \\ \hline
001 & sl & Shift left logical \\ \hline
010 & sra & Shift right arithmetic \\ \hline
011 & srl & Shift right logical \\ \hline
100 & ADD & A+B \\ \hline
101 & AND & A \& B \\ \hline
110 & OR & A $|$ B \\ \hline
111 & XOR & A \^{} B \\ \hline
\end{tabular}
\end{center}

The external interface of the ALU should be:

\textbf{Inputs}
\begin{itemize}
\item \texttt{InA}[15:0], \texttt{InB}[15:0] - Data input lines \texttt{InA} and \texttt{InB} (16 bits each).
\item \texttt{Cin} - A carry-in for the LSB of the adder.
\item \texttt{Oper}(2:0) – A 3-bit opcode that determines which operation should be performed and outputted.
The opcodes are shown in the table above.
\item \texttt{invA} - An invert-A input that causes the A input to be inverted before the operation is performed.
\texttt{invA} is active high, which means it inverts A when \texttt{invA} is 1.
\item \texttt{invB} - An invert-B input (also active high) that causes the \texttt{invB} input to be inverted before the
operation is performed.
\item \texttt{sign} – A flag to indicate whether signed-or-unsigned arithmetic should be performed in the ADD
function on the data lines (this also affects the \texttt{Ofl} output). The sign input is active high for signed
arithmetic and active low for unsigned arithmetic.
\end{itemize}

\textbf{Outputs}
\begin{itemize}
\item \texttt{Out}(15:0) – Output data from your ALU (16 bits).
\item \texttt{Ofl} - This should be high (1) if an overflow occurred (1 bit).
\item \texttt{Zero} - This should be high if the result is exactly zero (1 bit).
\end{itemize}

Other assumptions:
\begin{itemize}
\item You can assume 2's complement numbers.
\item In case of logic functions, Ofl is not asserted (i.e., kept logic low).
\end{itemize}

The top-level module definitions and a testbench are included in the download.  
\textbf{You must not change these top-level module names.}

Simulate and verify your design using the supplied testbench or create one yourself to test any of your
submodules. You must reuse the barrel shifter unit designed in Problem 1. \textbf{If/when you use files from
hw2\_1 in Problem 2, you must copy the files over from Problem 1 – you should not assume they are
in hw2\_2’s folder by default.}

In addition to the Verilog, you should also turn in schematics for each of your components. The schematics
may be handwritten or computer generated but must be legible if they are handwritten. The schematic files
should be named \textbf{schematic.pdf} and placed in the corresponding problems subdirectory (e.g.,
hw2\_1/schematic.pdf and hw2\_2/schematic.pdf). Although the schematics may seem simple for some of
these components, as your project gets bigger and bigger, you’ll find that drawing schematics of each
component and the bigger picture will make your task significantly easier.

\section*{Testing Instructions}

To run the testbench programs that we have distributed for this lab, use:
\begin{itemize}
\item \texttt{./run test <homework> [problem]}
\item Replace \texttt{<homework>} with the assignment name (e.g., \texttt{hw1}).
\item The \texttt{[problem]} argument is optional. If omitted, all tests for the selected homework assignment will be run.
\item For more detailed output, add \texttt{-v} (e.g., \texttt{./run test -v hw2 hw2\_1}).
\begin{itemize}
\item \textbf{Warning:} \texttt{-v} output can be extremely long in some contexts.
\end{itemize}
\end{itemize}

\section*{Verilog Rules}
Read the Verilog rules document at \texttt{assignments/tools/verilog\_rules/verilog\_rules.pdf} and adhere to the conventions.
Rules adherence will be checked when you run \texttt{./run test}.

\section*{Submission Instructions}

To build a submission for grading, run:
\begin{itemize}
\item \texttt{./run build\_submission <assignment>}
\item This will automatically run \texttt{./run test <assignment>}, then run the generated submission
through the integrated autograder.
\item A submission ZIP is created in \texttt{generated\_turnins/<assignment>/}.
\item If you run \texttt{build\_submission} multiple times for the same assignment, the generated zip files will be prepended with a number. The most recent submission should be the item with the highest number.
\item \textbf{IMPORTANT:} Running \texttt{./run build\_submission <assignment>} \textbf{DOES NOT} actively submit your assignment, it only generates the file you will submit. This file must be manually uploaded to Gradescope.
\end{itemize}

\end{document}
